% Related Literature — five-paragraph structure with paragraph headings,
% mirrored from v13 but rewritten to align with v15 single-thesis frame:
% construct/screen, cover-bidder population, descriptive scope, and
% complementarity-not-dominance against bid-distribution detectors.
% No reference to a Bajari-Ye section the v15 manuscript does not contain.

\section{Related Literature}
\label{sec:literature}

Four literatures bear on the screening problem. Each operates
under a constraint on data or on identification timing:
bid-coordination tests need the partition; proactive screens need
bid microdata; cover-bidding identification has been
retrospective; and the developing-economy literature has flagged
the resource gap but has not closed it. The construct's
positioning is at the conjunction of these constraints---not as a
more powerful test against any of them individually, but as a
screen that operates where the existing tests cannot, and that
integrates with them where they apply
(§\ref{sec:rob_imhof}).

\paragraph{Bid-coordination tests require the partition as input}
\citet{bajari2003deciding} establish the canonical econometric
framework for detecting collusion through exchangeability of bid
functions and conditional independence of bid residuals; the tests
achieve discriminating power against known cartels but require the
researcher to first partition bidders into suspected colluders and
a competitive fringe. The literature that builds on Bajari--Ye
inherits this constraint:
\citet{porter1993detection} and
\citet{porter1999ohio} detect bid rotation conditional on assumed
cartel membership; \citet{baldwin1997bidder} document collusion in
timber auctions identified through post-prosecution records;
\citet{conley2016detecting} infer bidder groups from co-bidding
patterns---the closest precedent in data parsimony, but recovering
group structure rather than individual flags;
\citet{clark2021collusion} and \citet{schurter2020identification}
study collusion dynamics in long-lived cartels with
ex-post-known cartel membership. The partition is constructed
\emph{ex post} or by assumption in every case. The frequent-loser
construct gives a candidate \emph{ex ante} firm-level flag
computable from contract-award records alone; bid-coordination
tests can then be applied with the construct's flag as input where
bid microdata are available.

\paragraph{Proactive bid-distribution screens require microdata
that procurement systems do not always preserve}
\citet{imhof2019detecting},
\citet{imhof2018screening},
\citet{huber2019machine}, and
\citet{wallimann2023machine} train machine-learning classifiers on
tender-level bid statistics (within-tender coefficient of
variation, kurtosis, skewness, normalized spread); the most
accurate implementations achieve AUC above $0.85$ against
ground-truth cartel cases in European procurement.
\citet{abrantesmetz2006a} pioneer the variance screen;
\citet{chassang2022robust} provide robust theoretical foundations
for moment-based detection; \citet{harrington2008detecting}
surveys the field. Every implementation operates at the
\emph{tender} level and requires per-bidder bid amounts. Many
real procurement systems do not preserve this layer at scale:
electronic platforms, including BEC during much of our sample,
archive winning prices and participant lists at the
contract-record layer but not bid-by-bid. The gap is operational,
not theoretical: what bid-distribution screens \emph{can} do with
full microdata is structurally larger than what oversight bodies
\emph{do} have. The frequent-loser construct runs on the second
data layer and integrates with the first
(§\ref{sec:rob_imhof}).

\paragraph{Cover-bidding identification in the empirical literature
has been retrospective} \citet{porter1999ohio},
\citet{pesendorfer2000study}, \citet{asker2010leniency}, and
\citet{marshall2012economics} document cover bidders in prosecuted
cartels through case records: the Ohio milk cartel, the Italian
highway cartel, the New York City stamp cartel, and a sequence of
cases reviewed by Marshall and Marx. \citet{athey2011comparing}
compare cover-bidding incentives across auction formats
theoretically. In every empirical case the cover bidders are
identified through leniency testimony, post-conviction records, or
judicial discovery---\emph{ex post}, after the cartel is
adjudicated. \emph{Ex ante} prospective identification of
cover-bidder candidates from public award-record data has not been
attempted in the empirical procurement literature, to our
knowledge. The construct inverts the retrospective workflow: by
flagging firms whose participation patterns are coherent with
cover bidding, the post-adjudication observation becomes a
pre-investigation screen.

\paragraph{Procurement institutions shape both the supply of
collusion and what screens can observe} A growing empirical
literature documents how administrative discretion, repeated
interaction, and political entrenchment alter procurement
outcomes outside the cartel-detection lens.
\citet{coviello2017tenure} show that longer political tenure
worsens procurement outcomes in Italy and is consistent with
tighter official--bidder relationships;
\citet{decarolis2020rules} show that discretionary procedures and
limits to competition jointly predict corruption risk in Italian
government contracting; \citet{decarolis2017corruption} documents
the institutional channel through which procurement-officer
characteristics moderate outcomes. These papers do not propose
participation-only screens, but they motivate the procuring-unit-size
oversight gradient we develop in §\ref{sec:results_oversight}:
where oversight infrastructure thins, the participation footprint
that the construct flags is amplified. The construct is most
naturally read as identifying environments where institutional
structure permits suspicious participation patterns, not as
recovering cartel membership firm by firm.

\paragraph{Procurement enforcement is hardest where data are
scarcest} The bid-rigging literature is dominated by
developed-economy settings with mature data infrastructure, in
which procurement spending represents $5$--$10\%$ of GDP. In
developing economies, where the share can exceed $\valPctTwelve$ of
GDP \citep{oecd2019government} and oversight resources are
thinner, settings are simultaneously more vulnerable to bid
rigging and less equipped to deploy bid-distribution detectors.
\citet{sanchezgraells2019screening} surveys the operational
limits of existing screens in low-resource jurisdictions. The
contribution generalizes in principle wherever a procurement
system records winners and participants electronically, regardless
of bid-by-bid archiving (§\ref{sec:lim_external}).

\paragraph{Synthesis} The four literatures share a common
structural constraint: each requires more information than typical
oversight bodies routinely have, or operates only after collusion
is established. Bid-coordination tests need the partition;
bid-distribution screens need microdata; cover-bidding
identification needs the post-adjudication record; and the
developing-economy literature has flagged the resource gap without
closing it. The frequent-loser construct addresses the conjunction.
It produces a candidate flag that bid-coordination tests can take
as input; it operates on the contract-award envelope that
bid-distribution screens cannot use; it shifts cover-bidding
identification from retrospective to prospective; and it is
deployable in low-resource settings. The contribution is not a
test that dominates any of the four benchmarks; it is a tool that
runs where the existing tests cannot and that integrates with them
where they apply.
